\documentclass[../../../../SoftwareRequirementsSpecification.tex]{subfiles}
\begin{document}
% Richiede le macro definite in src/macros/useCaseMacros.tex
% (incluse nel preambolo di SoftwareRequirementsSpecification.tex)

\ucstart
\begin{tabularx}{\textwidth}{|l|X|}
  \hline
  \ucid{A-03} \\ \hline
  \ucname{Richiesta Sessione di Allenamento} \\ \hline
  \ucactor{Atleta registrato} \\ \hline
  \ucdescription{L'atleta richiede una sessione di
    allenamento con il PT fornendo data e motivazione.} \\ \hline
  \ucprecond{L'atleta deve essere autenticato.} \\ \hline
  \ucpostcond{Viene creata una sessione di
    allenamento in attesa di conferma dal PT, il quale riceve
    una mail di notifica.} \\ \hline
  \ucmainflow{
    \item L'atleta atterra sulla pagina della lista sessioni prenotate.
    \item L'atleta preme il tasto "Richiedi Appuntamento".
    \item Il sistema mostra il form di richiesta appuntamento.
    \item L'atleta inserisce i dati necessari, in particolare data e ora
      della sessione e la motivazione.
    \item L'atleta preme il tasto "Invia Richiesta".
    \item Il sistema verifica i dati nel database.
    \item Il sistema salva nel database la richiesta di sessione.
    \item Il sistema invia al PT una mail per informarlo
      della nuova richiesta.
  } \\ \hline
  \ucaltflow{
    \textbf{6a. Richiesta in attesa già presente} \newline
    - Il sistema mostra un messaggio di errore ("C'è già una richiesta
    di sessione in attesa di risposta."). \newline
    - Il flusso riprende dal punto 1. \newline
  } \\ \hline
\end{tabularx}
\end{document}
